\section{What one pixel knows}\label{sec:one-pixel}

\begin{quote}
\textbf{In plain words.} A light sensor records six numbers at a point and a
gravitational pixel ten, although a wave of either kind is described by only
four. The extra numbers are not idle. From one reading of one pixel, with no
lens and no second pixel, the field tells which way a wave came from, how many
bodies are pulling, how fast the pixel itself is moving through the matter
around it, and where a moving body is now rather than where it was. Over time
it pulls apart several waves crossing it at once, and from how the field
departs from a perfect wave it tells how far away their source is. Light at a
point can do much less, for a reason that is a matter of algebra.
\end{quote}

A wave of light or of gravity reaching a point is described by four numbers: two
for its direction and two for the strengths of its two polarizations. A sensor
of light at a point records six, the electric and magnetic fields; a
gravitational pixel records ten, Eqs.~\eqref{eq:tidal} and~\eqref{eq:magnetic}.
This section shows what the six extra numbers buy. Most of it follows from
writing them as one object.

\paragraph{One complex matrix.} In vacuum $\tE$ and $\tB$ are real, symmetric and
trace-free. Together they form the complex symmetric trace-free matrix
\begin{equation}\label{eq:W}
\mathcal{W} = \tE + i\,\tB ,
\end{equation}
five complex numbers. Two facts make this the natural form.

First, observers transform $\mathcal{W}$ by rotations of three-dimensional space with
complex angles. A pixel turned by a rotation $R$ reads $R\,\mathcal{W} R^{T}$. A pixel moving
with rapidity $\eta$ along a unit vector $\hat u$ (speed $c\tanh\eta$) reads
$\Lambda\,\mathcal{W}\Lambda^{T}$ with $\Lambda = \exp(i\eta\,\hat u\cdot\mathbf{L})$, where
$(L_k)_{ij} = -\epsilon_{kij}$ generates rotations: a change of velocity is a
rotation by an imaginary angle. (The Lorentz group is isomorphic to the group of
complex rotations, $SO(3,\mathbb{C})$; for the curvature this is the content of
Petrov's and Penrose's classifications \cite{Petrov2000,Penrose1960}.) The
eigenvalues of $\mathcal{W}$, and with them
\begin{equation}\label{eq:invariants}
I = \mathrm{tr}\,\mathcal{W}^2, \qquad J = \mathrm{tr}\,\mathcal{W}^3,
\end{equation}
are therefore the same for every observer at that point, however turned and
however moving.

Second, a plane wave is a perfect square. For a wave travelling along $\hat k$,
let $\mathbf{e}_1, \mathbf{e}_2$ span the plane across it with
$\mathbf{e}_1\times\mathbf{e}_2 = \hat k$, and let
$\mathbf{m} = (\mathbf{e}_1 + i\mathbf{e}_2)/\sqrt2$, a complex vector with
$\mathbf{m}\cdot\mathbf{m} = 0$. Since $\tB_{ij} = \epsilon_{ipq}\hat k_p\tE_{qj}$ in
the wave, and $\hat k\times\mathbf{m} = -i\mathbf{m}$,
\begin{equation}\label{eq:square}
\mathcal{W} = 2(\tE_+ - i\tE_\times)\,\mathbf{m}\mathbf{m}^{T},
\end{equation}
where $\tE_+$ and $\tE_\times$ are the wave's two polarizations, so that a wave's
$\mathcal{W}$ has rank one. Conversely the direction of travel is read from $\mathbf{m}$
as $\hat k = \mathrm{Re}(i\,\mathbf{m}\times\mathbf{m}^*)/|\mathbf{m}|^2$, which for
one wave of any polarization is the gravitational counterpart of the Poynting
vector:
\begin{equation}\label{eq:one-wave}
\hat n_i = -\,\frac{\epsilon_{ijk}\,\tE_{jl}\,\tB_{kl}}{\tE_{ab}\,\tE_{ab}} ,
\end{equation}
with $\hat n = -\hat k$ pointing from the pixel back to the source. The same
construction for light gives $\mathbf{E} + ic\mathbf{B} = a\,\mathbf{m}$, a vector
rather than a square; the difference between the two, a square against a single
power, is the difference between spin two and spin one, and it is the source of
almost everything below. Matter neither reflects nor absorbs gravitational waves
(Section~\ref{sec:no-lens}), so none of what follows is spoiled by echoes, as
direction finding with radio or sound is indoors.

\paragraph{Four stars.} Every reading has four directions of its own. Writing a
complex null vector by its stereographic coordinate $z$ on the sky,
$\mathbf{m}(z) = (1 - z^2,\, i(1 + z^2),\, -2z)$, the quartic equation
\begin{equation}\label{eq:stars}
\mathbf{m}(z)^{T}\,\mathcal{W}\,\mathbf{m}(z) = 0
\end{equation}
has four roots, four points on the sphere of directions around the pixel: the
principal null directions of the field \cite{Penrose1960}, which we call its
stars. For one plane wave all four coincide on its direction of travel. For one
body at rest they fall in two pairs, on the line to the body and opposite it.
Where two or more bodies pull, they are four distinct points. The same
distinctions can be read without solving anything, from the invariants: a
repeated eigenvalue, the mark of a single body, makes
$1 - 6J^2/I^3$ vanish, and a single wave makes $I = J = 0$ (Petrov's types D and
N). One reading therefore tells whether the field could be that of one body,
whether it is a pure wave, or whether it needs more. In the reference
implementation the quantity $1 - 6J^2/I^3$ is below $10^{-8}$, rounding, for a
single moving body, and $0.02$ for two bodies at comparable distances.
Light at a point has two stars, not four: its $\mathbf{E} + ic\mathbf{B}$ has
only one complex invariant.

\paragraph{Ten numbers: four invariants, three angles, three velocities.} A
reading whose eigenvalues are distinct can be written
\begin{equation}\label{eq:polar}
\mathcal{W} = P\,\mathrm{diag}(w_1, w_2, w_3)\,P^{T}, \qquad
P = R_0\,e^{iK_0},
\end{equation}
with $P$ a complex rotation ($P^TP = 1$), split into a real rotation $R_0$ and an
imaginary one. The three complex eigenvalues, which sum to zero, are four
numbers no observer can change. The rotation $R_0$ gives three angles, the
orientation of the field's own axes. The imaginary part,
$e^{2iK_0} = P^{\dagger}P$, gives three more: the velocity of the one frame in
which $\tE$ and $\tB$ share their axes, the field's own rest frame. Ten numbers,
read as four, three and three.

The consequences are immediate. For a scene of masses at rest, $\tB = 0$ in the
scene's frame, so its eigenvalues are real; a pixel moving through it at any
speed still finds real eigenvalues, and finds in $P$ the velocity of the scene
relative to itself. The pixel is a speedometer: it reads its own velocity through
an unknown room, at any speed, with no map of the room. In the reference
implementation, for a pixel moving at $0.6\,c$ past four unknown masses, the
rapidity is recovered to $4\times10^{-8}$. To first order in the speed, in the
principal axes of $\tE$ with eigenvalues $\lambda_a$,
\begin{equation}\label{eq:speedometer}
\tB_{bc} = \frac{v_a}{c}\,(\lambda_c - \lambda_b) \quad ((a,b,c)\ \text{cyclic}),
\qquad \tB_{aa} = 0 ,
\end{equation}
so a pixel whose readings carry white noise $\sigma$ on each component reads
each component of its velocity to within
$c\,\sigma/\sqrt2\,|\lambda_b - \lambda_c|$. A pixel that reads the tide of a
room, $10^{-8}\,\mathrm{s^{-2}}$ near a person, to one part in $10^{8}$, reads its
own speed to about two metres per second.

The converse is an alarm. A pixel moving through a static scene, however fast,
always finds real eigenvalues. If they are complex, something in the scene
moves relative to the rest of it, or turns: the pixel's own motion can never
set it off. The phase of the eigenvalues measures turning. Outside a spinning
body of mass $M$ and angular momentum $J_s = Mca$, the eigenvalues are
proportional to $(r - ia\cos\vartheta)^{-3}$, as in Kerr's solution
\cite{Kerr1963}, so their phase, $3a\cos\vartheta/r$ for $a \ll r$, reads the
spin along the line of sight; for the Earth seen from above a pole it is
$1.5\times10^{-6}$.

Light cannot do this. Its field at a point has only two invariants,
$\mathbf{E}^2 - c^2\mathbf{B}^2$ and $\mathbf{E}\cdot\mathbf{B}$, and in the frame
where $\mathbf{E}$ is parallel to $\mathbf{B}$ a boost along their common axis
and a turn about it change nothing. A sensor of light is at best a
two-dimensional speedometer. A single body does the same to gravity's field
(type D): coming straight at the pixel, even at $0.3\,c$, it shows no $\tB$ at
all, and its radial speed is unreadable. Two bodies, or one that is not round,
remove the blind spot.

\paragraph{A moving body, exactly.} The four stars of a single body moving at any
speed are its two pairs, displaced by aberration toward its direction of motion.
Their midpoint, as a vector, is exactly its velocity across the line of sight:
\begin{equation}\label{eq:star-velocity}
\tfrac12\,(\hat p_+ + \hat p_-) = \frac{\mathbf{v}_\perp}{c} ,
\end{equation}
with $\hat p_\pm$ the directions of the two pairs and $\mathbf{v}_\perp$ the
velocity across the line from the body to the pixel. The reference
implementation finds the two sides equal to $2\times10^{-7}$ at $0.9\,c$. The
ratio $|\tB|/|\tE| = \sqrt3\,v_\perp/c$ of Section~\ref{sec:separation} is its
first-order form. A body coming straight on shows no sideways velocity and
cannot be ranged (Section~\ref{sec:separation}), but because its tide grows as
$r^{-3}$, the pixel reads when it will arrive,
\begin{equation}\label{eq:contact}
-\frac{r}{\dot r} = \frac{3\,\tE_\parallel}{\dot{\tE}_\parallel},
\qquad
\frac{\tE_\parallel\,\ddot{\tE}_\parallel}{\dot{\tE}_\parallel^2} = \frac43 - \frac{r\ddot r}{3\dot r^2},
\end{equation}
and from the second ratio whether it is coasting ($4/3$), braking or driven.

\paragraph{Separating waves at one point.} Several waves crossing a pixel at
once add their matrices, and the algebra of squares lets one pixel pull them
apart.

At one instant, two waves give $\mathcal{W} = c_1\mathbf{m}_1\mathbf{m}_1^{T} +
c_2\mathbf{m}_2\mathbf{m}_2^{T}$, of rank two, so $\det\mathcal{W} = 0$: two real
conditions that tell two crossing waves from the pull of a mass, whose
$\det\mathcal{W} = -2(GM/r^3)^3$. The vector $\mathbf{v}$ with $\mathcal{W}\mathbf{v} = 0$ is
$\mathbf{m}_1\times\mathbf{m}_2$; the null vectors orthogonal to it are the
roots of one quadratic, and they are $\mathbf{m}_1$ and $\mathbf{m}_2$, which give
both directions and then both amplitudes. This is Sylvester's theorem on sums of
two fourth powers \cite{Sylvester1851} in the language of the sky. Light cannot
do it at all: a single wave already uses four of its six numbers.

At one frequency, three waves can be separated. Writing the reading as
$\mathcal{W}(t) = \mathcal{W}_+e^{-i\omega t} + \mathcal{W}_-e^{i\omega t}$, both parts are sums
over the same three squares, $\mathcal{W}_\pm = \sum_a \alpha_a^\pm\,
\mathbf{m}_a\mathbf{m}_a^{T}$, and
\begin{equation}\label{eq:pencil}
\det(\mathcal{W}_+ - \lambda\,\mathcal{W}_-) = \det[\mathbf{m}_1,\mathbf{m}_2,\mathbf{m}_3]^2
\prod_{a=1}^{3}(\alpha^+_a - \lambda\,\alpha^-_a) .
\end{equation}
Its three roots give the ratios $\lambda_a = \alpha_a^+/\alpha_a^-$; at each, one
wave drops out, the null vector of what remains is the cross product of the other
two $\mathbf{m}$, and cross products of those give all three directions. For a
circular binary seen at inclination $\iota$, $|\lambda| = \tan^{\pm4}(\iota/2)$,
so the roots also give the three inclinations. Three binaries in one frequency
bin, which a single interferometer cannot tell apart, are three stars to one
pixel (to $10^{-9}$ in the reference implementation).

Over time, four. If the waves come from fixed directions but carry any
waveforms at all, the readings span a subspace of the five-dimensional space of
matrices; for four waves it is four-dimensional, and there is one matrix $A$ with
$\mathrm{tr}(A\mathcal{W}) = 0$ for every reading. Each wave's $\mathbf{m}$ then lies on
two conics,
\begin{equation}\label{eq:four-waves}
\mathbf{m}^{T}\!A\,\mathbf{m} = 0, \qquad \mathbf{m}\cdot\mathbf{m} = 0 ,
\end{equation}
and two conics meet in four points: the four directions. For light the same
argument gives two. In general a point sensor of a field of spin $s$ separates
$s$ plane waves at one instant, $s + 1$ at one frequency and $2s$ from a time
history; for light, $1$, $2$ and $2$, the known limit for electromagnetic vector
sensors \cite{NehoraiPaldi1994,HochwaldNehorai1996}; for gravity, $2$, $3$ and
$4$. In the reference implementation four waves with random waveforms are
located to $10^{-14}$, and to $10^{-4}$ radian with noise of $10^{-4}$ of the
signal on every component over three hundred readings.

Uncorrelated sources can be separated from second moments alone. In the
frequency domain the two helicities occupy separate halves of the ten numbers,
and the brightness and the circular polarization of the sources give two
five-by-five matrices built from the same five steering vectors; when the
sources' circular fractions differ, diagonalizing one in the frame that whitens
the other separates five sources in closed form, with their directions,
brightnesses and circular fractions (checked in the reference implementation).

\paragraph{What one pixel sees of the sky.} For a sky full of independent
sources overlapping in frequency, what a pixel can learn at each frequency is
its ten-by-ten spectral matrix, one hundred real numbers. Split by helicity, it
holds three blocks: the brightness $I$ and the circular polarization $V$ of the
sky in two five-by-five Hermitian blocks, and the linear polarization in a
complex symmetric one. The brightness enters through the kernel
\begin{equation}\label{eq:pixel-kernel}
\begin{gathered}
\sum_{A=+,\times} e^{A}_{ij}\, e^{A}_{kl}
= P_{ik}P_{jl} + P_{il}P_{jk} - P_{ij}P_{kl},\\
P_{ij} = \delta_{ij} - \hat n_i \hat n_j ,
\end{gathered}
\end{equation}
a polynomial of degree four in the components of $\hat n$, so the pixel sees the
spherical harmonics of $I$ of degree four and below, twenty-five numbers; it
sees as many of $V$, and eighteen of the linear polarization, which carries spin
weight four and enters only at degree four. That is
\begin{equation}\label{eq:sky-count}
25 + 25 + 18 = 68
\end{equation}
numbers of the sky, all independent; the other thirty-two of the hundred vanish
for any sky of gravitational waves and are checks against sensor noise or
anything else. The brightness and the circular polarization never mix at one
pixel: one lives in the real part of the spectral matrix, the other in its
imaginary part. (They mix in the cross-correlation of two separated detectors
\cite{AllenOttewill1997,SetoCooray2004}, through the phase between them.) The
same count for a point sensor of spin $s$ gives $2(2s+1)^2 + 2(4s+1)$: $28$ of
$36$ for light. For sound, a pressure and a velocity, it is $9$ of $16$. All
three counts are checked in the reference implementation.

One of the sixty-eight numbers is the handedness of the whole sky. An isotropic
sky of brightness $I$ and net circular polarization $V$ puts exactly
$(V/I)\langle\tE\cdot\tE^*\rangle$ into the imaginary part of
$\langle\tE\cdot\tB^*\rangle$, so one pixel measures
whether the sky is chiral, which a single planar interferometer cannot
\cite{SetoTaruya2007}.

The rest of the sky is invisible to second moments. The twelve corners of a
regular icosahedron form a spherical design of degree five: they average every
polynomial of degree five or less exactly as the whole sphere does
\cite{Delsarte1977}. Twelve equal, unpolarized, independent sources at those
corners therefore give a pixel exactly the second moments of a uniform sky,
while six at the corners of an octahedron, a design of degree three only, do
not; the reference implementation finds the first pair equal to one part in
$10^{15}$. Sources that are not Gaussian say more: the fourth moments of the
pixel are kernels of degree eight, which no icosahedron averages. Two pixels a
wavelength apart tell the two skies apart as well, because their
cross-correlation carries the phase $e^{i\omega\hat n\cdot\Delta\mathbf x/c}$,
which no finite design averages exactly.

\paragraph{How far away is a wave source?} A single plane wave has no distance.
The field of a real source at distance $r$ is not quite a plane wave: it carries
terms falling as $1/r^2$, $1/r^3$ and beyond, organized by how they turn about
the line of sight \cite{Sachs1962}, and a pixel sees them. At one instant the
departure of the oscillating part of $\mathcal{W}$ from rank one, the ratio of its
second singular value to its first, is
\begin{equation}\label{eq:ranging-instant}
\frac{s_2}{s_1} = \frac{C(\iota, \theta)}{(kr)^2} + O\bigl((kr)^{-3}\bigr),
\end{equation}
not $1/kr$: the first correction only tilts the apparent $\mathbf{m}$, which at one
instant cannot be told from a tilt of the source. For two masses turning in a
circle, seen at inclination $\iota$ and wave phase $\theta$,
\begin{equation}\label{eq:ranging-C}
\begin{gathered}
C = \frac{2\sin^2\!\iota\,\sqrt{A^2 + B^2}}{D}, \qquad
B = 2\cos\iota\,\sin\theta\cos\theta,\\
A = (1 - 3\cos^2\!\iota)\cos^2\!\theta - 4\cos^2\!\iota\,\sin^2\!\theta,\\
D = (1 + \cos^2\!\iota)^2\cos^2\!\theta + 4\cos^2\!\iota\,\sin^2\!\theta,
\end{gathered}
\end{equation}
which is $2$ for any phase seen edge on and vanishes at the angle
$\cos^2\iota = 1/3$. Over a cycle more can be done. The two frequency parts
$\mathcal{W}_\pm$ depart from rank one by $2\tan^{\pm2}(\iota/2)/(kr)^2$, so that
$kr = [2\sqrt{s_1(\mathcal{W}_+)s_1(\mathcal{W}_-)}/s_2(\mathcal{W}_+)]^{1/2}$ needs no knowledge of the
inclination. And the first-order tilt, invisible at one instant, differs between
the two helicities: the directions read from the two parts differ by
\begin{equation}\label{eq:helicity-split}
\Delta\theta_\pm = \frac{8\,|\cot\iota|}{kr} .
\end{equation}
A difference of directions measured at signal-to-noise ratio $\rho$ ranges the
source out to $kr$ of order $\rho$, where the rank departure reaches only $kr$
of order $\sqrt\rho$. A source seen at $\rho = 10^{3}$ can be ranged out to
several hundred reduced wavelengths: at $100\,\mathrm{Hz}$, a few hundred
thousand kilometres. The
reference implementation checks Eqs.~\eqref{eq:ranging-C}
and~\eqref{eq:helicity-split} against the exact field of a turning pair to
within a few per cent at $kr = 80$ and $0.5$ per cent at $kr = 200$.

\paragraph{The camera sees the present.} The field of a body moving uniformly is
its static field carried along, so it depends on where the body is now. The
stretch axis of $\tE$ points at the body's present position, at any speed (to
$10^{-7}$ radian at $0.3\,c$ in the reference implementation), while its light
shows its position a light-travel time ago, displaced by the aberration angle,
about $v_\perp/c$. Near two masses in orbit the changing part of the field is
Newton's field of their present positions to order $(kr)^2$, while their
retarded positions would be off by $kr$ itself: at $kr = 0.1$, a tenth of a per
cent against ten per cent. The first-order retardation cancels because mass and
momentum are conserved. The gravitational Sun is therefore $20.5''$ ahead of the
optical Sun along its yearly path across the sky. Laplace took the absence of
gravitational aberration to mean that gravity is far faster than light; it is
what general relativity predicts \cite{Carlip2000}. A camera that sees mass sees
it on its own slice of simultaneity. When a body changes its motion, the field
goes on pointing at where it would have been until the news, travelling at the
speed of light as a wave, arrives.

\paragraph{Beyond general relativity, in one reading.} Other theories of gravity
allow up to six polarizations, and waves that travel at other speeds
\cite{EardleyLeeLightman1973}. For any plane wave of the metric, of any
polarization and phase speed $v_p$, the second Bianchi identity gives, with no
field equations assumed,
\begin{equation}\label{eq:bianchi-wave}
\tB_{ij} = \frac{c}{v_p}\,\epsilon_{ipq}\hat k_p\,\tE_{qj} .
\end{equation}
So one reading of $\tE$ (now with a trace, six numbers) and $\tB$ (with an
antisymmetric part, eight) gives the direction ($\hat k$ is the null vector of
$\tB$ from the left), the phase speed, all six polarizations, and five checks that
they belong to one wave. In general relativity $v_p = c$ and only the two tensor
polarizations appear. A massive graviton and a massive scalar differ even in
sign: with $n = c/v_p$, the helicity-zero wave of a Fierz--Pauli graviton
\cite{FierzPauli1939} has longitudinal and breathing amplitudes in the ratio
$-2(1 - n^2)$, a scalar wave $+(1 - n^2)$. These are checked against curvature
computed from the metric in the reference implementation. Testing general
relativity this way at the level of present bounds would need extraordinary
signal-to-noise ratios; what is new is that it needs only one point.

\paragraph{How precisely.} With white noise $\sigma$ on each of the ten numbers
and $\rho = |\mathcal{W}|/\sigma$, one pixel places one wave to within $1/\rho$ radian in
each direction, against $\sqrt2/\rho$ for a sensor of light at the same
signal-to-noise ratio; it reads its own velocity as given after
Eq.~\eqref{eq:speedometer}; and it ranges a moving mass with the precision that
Section~\ref{sec:separation} gives.

\paragraph{Matter at the pixel.} Inside matter the pixel reads the density
locally, whatever the matter is: in full,
$\mathrm{tr}\,\tE = 4\pi G(\rho + 3p/c^2)$, less $2\Omega^2$ if the pixel's frame
turns at $\Omega$. The dark matter around the Sun, $0.4\,\mathrm{GeV\,cm^{-3}}$,
gives $6\times10^{-31}\,\mathrm{s^{-2}}$, by its mass alone and with no other
interaction. Between two plates a micrometre apart the Casimir energy is
negative and its pressure large, and the trace is negative,
$-8\times10^{-30}\,\mathrm{s^{-2}}$. The cosmological constant gives
$-\Lambda c^2 = -1\times10^{-35}\,\mathrm{s^{-2}}$. Where radiation streams
through, as neutrinos from a supernova do, the trace gives its energy flux and the
antisymmetric part of $\tB$ its direction (Section~\ref{sec:pixel}), so that
one reading inside the passing shell points back to the star, even when the
collapse is perfectly spherical and sends out no waves at all.

\begin{table}
\centering
\begin{tabular}{p{0.46\linewidth}p{0.12\linewidth}p{0.13\linewidth}p{0.14\linewidth}}
At one point & Sound & Light & Gravity \\ \hline
Numbers recorded & 4 & 6 & 10 \\
Numbers no observer can change & -- & 2 & 4 \\
Own velocity readable, components & -- & 2 & 3 \\
Waves separated at one instant & 1 & 1 & 2 \\
Waves separated at one frequency & -- & 2 & 3 \\
Waves separated over time & 2 & 2 & 4 \\
Sky numbers per frequency (of) & 9 (16) & 28 (36) & 68 (100) \\
Direction of one wave, standard deviation & -- & $\sqrt2/\rho$ & $1/\rho$ \\
\end{tabular}
\caption{What a sensor at one point can know, for sound (pressure and velocity),
light (electric and magnetic fields) and gravity (the ten numbers of a pixel).
Separation counts are for plane waves from distinct directions; over time the
waveforms are arbitrary. A dash marks an entry that does not
apply.}\label{tab:one-pixel}
\end{table}
